\chapter{Trabajo relacionado y Estado del Arte (revisar)} \label{chp:state-of-the-art}

Este capítulo analiza el contexto tecnológico y normativo en el que se enmarca este Trabajo de Fin de Grado \cite{primera_entrega} [1]. En primer lugar, se expone la situación actual de las comunicaciones con la Seguridad Social mediante el Sistema RED y la falta de interfaces modernas \cite{blog_diario} [2, 3]. A continuación, se evalúan las soluciones alternativas disponibles en el mercado, como las APIs de terceros \cite{blog_diario} [4, 5]. Posteriormente, se describe el estado del arte de las tecnologías de desarrollo ágil empleadas, específicamente el paradigma \textit{Low-Code} y los servicios \textit{Backend-as-a-Service} (BaaS) \cite{documento_diseno, proyecto_alta} [6, 7]. Finalmente, se aborda la literatura relacionada con las arquitecturas \textit{Multi-Tenant} y el estricto cumplimiento del Reglamento General de Protección de Datos (RGPD) \cite{primera_entrega, proyecto_alta} [1, 8, 9].

%%%%%%%%%%%%%%%%%%%%%%%%%%%%%%%%%%%%%%%%%%%%%%%%%%%%%%%%%%%%%%%%%%%%%%%%%%%%%%%%
%%%%%%%%%%%%%%%%%%%%%%%%%%%%%%%%%%%%%%%%%%%%%%%%%%%%%%%%%%%%%%%%%%%%%%%%%%%%%%%%

\section{El Sistema RED y la comunicación oficial con la TGSS} \label{sct:state:sistema_red}

El Sistema RED (Remisión Electrónica de Datos) es el canal de comunicación obligatorio establecido por la administración pública española para que las empresas realicen trámites de afiliación, como altas, bajas y variaciones de datos, así como la cotización a la Seguridad Social \cite{blog_diario} [2]. En el estado actual de la tecnología gubernamental, la Tesorería General de la Seguridad Social (TGSS) no dispone de una API pública web sencilla y directa, basada en estándares modernos como REST/JSON, que pueda ser consumida fácilmente desde plataformas en la nube \cite{blog_diario} [3].

Tradicionalmente, esta comunicación requiere el uso de programas de escritorio específicos proporcionados por la administración, como la aplicación SILTRA o RED Directo \cite{blog_diario} [2]. Estas herramientas exigen a los despachos laborales y gestorías manejar ficheros físicos, obtener autorizaciones previas como usuarios RED y, de forma imprescindible, gestionar certificados digitales locales para cada cliente de manera individual \cite{blog_diario} [3]. Este enfoque representa un importante cuello de botella tecnológico, ya que dificulta la integración directa desde aplicaciones web o móviles modernas, exigiendo habitualmente el desarrollo de un \textit{software} de comunicaciones especializado o la implantación de un \textit{middleware} complejo \cite{blog_diario} [3].

%%%%%%%%%%%%%%%%%%%%%%%%%%%%%%%%%%%%%%%%%%%%%%%%%%%%%%%%%%%%%%%%%%%%%%%%%%%%%%%%
%%%%%%%%%%%%%%%%%%%%%%%%%%%%%%%%%%%%%%%%%%%%%%%%%%%%%%%%%%%%%%%%%%%%%%%%%%%%%%%%

\section{Soluciones alternativas y APIs de terceros} \label{sct:state:apis_terceros}

Dada la extrema dificultad técnica de integrar el Sistema RED directamente desde plataformas \textit{frontend} o en la nube, el mercado ha evolucionado hacia la provisión de soluciones intermedias \cite{blog_diario} [3, 4]. Existen empresas proveedoras de \textit{software} de Recursos Humanos y nóminas que ofrecen servicios de pasarela y abstracción \cite{blog_diario} [4]. Estas empresas exponen una API moderna (REST/JSON) orientada a desarrolladores, que se encarga de asimilar toda la complejidad del Sistema RED, manejando por detrás los certificados, los formatos de ficheros y los protocolos rígidos de la administración \cite{blog_diario} [4, 5]. Un ejemplo destacado de este tipo de soluciones es la pasarela proporcionada por la empresa \textit{CreateApps} \cite{blog_diario} [10, 11].

El uso de estas APIs de terceros permite una integración sencilla con servicios en la nube (como Firebase Cloud Functions) mediante llamadas HTTP estándar \cite{blog_diario} [5]. Sin embargo, presentan como desventaja una dependencia tecnológica vital hacia un proveedor de terceros para una funcionalidad crítica y un coste económico por transacción o suscripción \cite{blog_diario} [5]. Debido a estas limitaciones y a los costes asociados, en el ámbito académico y de prototipado del presente trabajo, se ha optado por el diseño y desarrollo de una \textit{Mock API} (API simulada) desplegada de forma local \cite{blog_diario, primera_entrega} [1, 11, 12]. Este componente emula las respuestas y validaciones de la administración pública para poder evaluar la robustez del sistema y el flujo funcional completo sin incurrir en dependencias ni gastos \cite{blog_diario} [11, 12].

%%%%%%%%%%%%%%%%%%%%%%%%%%%%%%%%%%%%%%%%%%%%%%%%%%%%%%%%%%%%%%%%%%%%%%%%%%%%%%%%
%%%%%%%%%%%%%%%%%%%%%%%%%%%%%%%%%%%%%%%%%%%%%%%%%%%%%%%%%%%%%%%%%%%%%%%%%%%%%%%%

\section{Tecnologías Low-Code y Backend-as-a-Service (BaaS)} \label{sct:state:lowcode_baas}

En los últimos años, el desarrollo de aplicaciones empresariales ha adoptado enfoques orientados a reducir drásticamente el tiempo de comercialización (\textit{Time to Market}) y la complejidad de la gestión de la infraestructura subyacente. El paradigma \textit{Low-Code} permite la creación rápida de interfaces de usuario multiplataforma mediante un modelado visual \cite{documento_diseno} [6]. En este contexto, \textit{FlutterFlow} destaca como un entorno de desarrollo que genera código nativo estructurado basado en el \textit{framework} Flutter, permitiendo construir el diseño UI/UX de formularios complejos y habilitar validaciones algorítmicas críticas en tiempo real de manera muy eficiente \cite{documento_diseno, primera_entrega} [1, 6].

De forma complementaria, el estado del arte en la gestión de infraestructura en la nube se apoya en arquitecturas \textit{Backend-as-a-Service} (BaaS) como Firebase, impulsado por Google \cite{proyecto_alta} [7, 13]. Esta plataforma abstrae la administración de servidores dedicados y provee un ecosistema de servicios listos para su uso: Firebase Authentication para la gestión segura de accesos mediante roles; Cloud Firestore como base de datos NoSQL documental orientada a la escalabilidad en tiempo real; y Cloud Functions \cite{documento_diseno, proyecto_alta} [6, 7]. Estas funciones en la nube (ejecutadas típicamente bajo Node.js) actúan como un \textit{middleware} seguro, permitiendo centralizar la lógica de negocio compleja y las integraciones con APIs externas (como el SEPE o la TGSS) en un entorno protegido, evitando por completo la exposición de credenciales o secretos criptográficos en el lado del cliente \cite{documento_diseno, proyecto_alta} [6, 7, 14].

%%%%%%%%%%%%%%%%%%%%%%%%%%%%%%%%%%%%%%%%%%%%%%%%%%%%%%%%%%%%%%%%%%%%%%%%%%%%%%%%
%%%%%%%%%%%%%%%%%%%%%%%%%%%%%%%%%%%%%%%%%%%%%%%%%%%%%%%%%%%%%%%%%%%%%%%%%%%%%%%%

\section{Arquitecturas Multi-Tenant y Cumplimiento del RGPD} \label{sct:state:multitenant_rgpd}

Para dar servicio a gestorías de Recursos Humanos que administran de forma simultánea a múltiples empresas cliente, los sistemas de información modernos deben implementar de base una arquitectura \textit{Multi-Tenant} (multi-inquilino) \cite{primera_entrega} [1]. En bases de datos de naturaleza NoSQL como Firestore, este modelo arquitectónico se resuelve utilizando patrones de referencias cruzadas entre documentos (\textit{Document Reference}), asociando de manera inequívoca y relacional a cada empleado y cada usuario de la gestoría con la empresa o empresas a las que pertenecen \cite{documento_diseno} [15, 16].

Este aislamiento y segregación estricta de la información es un requerimiento técnico innegociable para el cumplimiento íntegro del Reglamento General de Protección de Datos (RGPD) en el ámbito europeo \cite{primera_entrega, proyecto_alta} [1, 9]. Al tratar diariamente con datos personales, fiscales y contractuales de alta sensibilidad (DNI/NIE, cuentas bancarias IBAN, afiliación a la Seguridad Social, contratos laborales), el sistema debe garantizar controles y mecanismos granulares de autorización \cite{proyecto_alta} [8]. De este modo, un administrativo o usuario de Recursos Humanos únicamente podrá generar registros de alta o visualizar el historial de las compañías a las que haya sido explícitamente autorizado por un administrador de sistema \cite{blog_diario, proyecto_alta} [8, 17]. Adicionalmente, el RGPD y los marcos de buenas prácticas de auditoría exigen mantener una trazabilidad ininterrumpida de las operaciones sobre los datos, por lo que este tipo de arquitecturas incorporan de diseño colecciones históricas (tales como \textit{altas\_history}), encargadas de auditar de forma inmutable cada trámite, identificando el responsable de la acción, la empresa afectada, el estado del envío y la fecha de tramitación (sello temporal) \cite{blog_diario, proyecto_alta} [18, 19].